\chapter{Reproduction of the Test Cases of Klempt et al.\ 2024}
\label{app:klempt2024}

% Written 3 Oct 2026 from KLEMPT2024_REPRODUCTION.md (sec. 9-15) and the
% scripts named there (JAXFEM/klempt2024_*.py). Numbers as in those files.

Before coupling the field model of Klempt et al.~\cite{Klempt2024DiffusionDrivenGrowth}
into the finite element framework, I solved its two test cases independently
in Python, to see whether the equations and parameters as printed give the
published results. They do not in every respect, and this appendix records
what is needed to reproduce them. None of the results of
Chapter~\ref{ch:ansys} depends on this reproduction; it informs the
background of Section~\ref{sec:theory_field} and the questions to the authors.

\section{Set-up}

The domain is the paper's cube of $20~\mu$m edge length on a regular grid of
$21^3$ nodes. Equations~\eqref{eq:theory_phi}--\eqref{eq:theory_alpha} are
integrated with explicit time steps and upwinded gradients in the front term;
$\phi$ is limited to $[0,1]$. All parameters are those of
Table~\ref{tab:theory_klempt2024}. The two test cases are:
\begin{itemize}
  \item \emph{4.1}: a spherical seed of $5~\mu$m radius in the centre, the nutrient held at
        $c=1$ in one corner of the cube. Read from the paper's Fig.~3 and
        Table~3, the source is a block of about $2~\mu$m at the corner; the
        strip of the schematic Fig.~2 fits worse.
  \item \emph{4.2}: a disk of $5~\mu$m diameter just above a bottom face
        held at $c=1$, with the consumption $g=10^8$ (``high'' nutrient level)
        and $g=10^{10}$ (``low'').
\end{itemize}
The published curves of the mean $\phi$ and mean $c$ (Figs.~4 and~7) were
read off the PDF by colour and compared with the reproduction by their RMS
difference over $T^*\in[0,1]$.

\section{The printed equations}

Solved as printed, the front term of Eq.~\eqref{eq:theory_phi} moves the seed
towards the nutrient rather than enlarging it: the factor
$\mathbf n_{\nabla\phi}\cdot\mathbf n_{\nabla c}$ is positive on the face of
the colony that looks towards the nutrient and negative on the other, so one
face grows while the other recedes. The mean $\phi$ of Fig.~4 is not reached
with any numerical variant tried (seed held fixed, no limiter, a second-order
upwind scheme), nor with any time scale. Two further observations point the
same way: in case 4.2 ``high'' the front advances at about $150~\mu$m per
$T^*$, three times the largest speed $r\,c/(k+c)$ that Table~2 allows, and in
Table~4 of the paper the biofilm grows away from the nutrient face, while the
printed term drives it towards it.

The consumption term gives a clean check. With the consumption
$g\,\phi\,c$, proportional to the nutrient, the steady nutrient level of
case 4.2 ``high'' is $\tanh(2)/2=0.482$, against $0.489$ in the paper. The
printed form $g\,\phi$ gives $0.139$.

\section{What reproduces the curves}

The following changes, each stated as an assumption, reproduce all three
curves without changing a value of Table~2:
\begin{enumerate}
  \item growth on every face of the colony, with the front term
        $\|\nabla_X\phi\|\,\dfrac{r\,c}{k+c}\bigl[w+(1-w)\,|\mathbf n_{\nabla\phi}\cdot\mathbf n_{\nabla c}|\bigr]$
        and an isotropic share $w=0.5$ (the text of the paper says that growth
        perpendicular to the nutrient gradient is smaller, not zero);
  \item consumption $g\,\phi\,c$;
  \item a time scale $s$ per simulation, applied to $r$, $\beta$ and
        $k_\alpha$: the paper normalises time by a reference time ``chosen
        based on the conditions for growth at hand'', which is not given.
\end{enumerate}
Table~\ref{tab:app_klempt_rms} lists the resulting differences. For $w$
between $0.35$ and $0.7$ the worst of the three stays between $0.029$ and
$0.032$, so the curves do not determine $w$ precisely.

\begin{table}[htbp]
  \centering
  \caption[Reproduction of Klempt et al.\ 2024, mean curves]{RMS difference
    between the reproduction and the curves of Klempt et al.\ 2024 (mean
    $\phi$ and mean $c$, averaged), with the changes listed in the text.}
  \label{tab:app_klempt_rms}
  \begin{tabular}{lccc}
    \hline
    test case & time scale $s$ & RMS, $w=0.5$ & RMS, $w=0$\\
    \hline
    4.1 (Fig.~4) & $1$ & $0.013$ & $0.022$\\
    4.2 high (Fig.~7) & $5$ & $0.021$ & $0.038$\\
    4.2 low (Fig.~7) & $3$ & $0.030$ & $0.07$\\
    \hline
  \end{tabular}
\end{table}

The shapes agree in kind (Figures~\ref{fig:app_klempt_fig3}
and~\ref{fig:app_klempt_table4}). In case 4.1 the colony becomes an egg drawn
out towards the nutrient corner and reaches it, as in the paper. In case 4.2
the early dome and the final layer of the ``low'' case look like the paper's;
in the ``high'' case the reproduction spreads to the side walls earlier, where
the paper's colony rises with straight sides. In the variant without the
isotropic share ($w=0$), varying the regularisation $\beta$ and the upwind
scheme did not change the shapes, so the remaining
differences lie in the equations or the set-up rather than in the numerics.

\begin{figure}[htbp]
  \centering
  \includegraphics[width=\linewidth]{\ansysfigdir fig_klempt2024_fig3_blend.png}
  \caption[Test case 4.1, reproduction and paper]{Test case 4.1: isosurface
    $\phi=0.8$ at $5$, $25$, $45$ and $100\,\%$ of the simulated time. Top:
    reproduction with the changes listed in the text ($w=0.5$, $s=1$).
    Bottom: Fig.~3 of Klempt et al.\ 2024 (CC BY 4.0).}
  \label{fig:app_klempt_fig3}
\end{figure}

\begin{figure}[htbp]
  \centering
  \includegraphics[width=0.85\linewidth]{\ansysfigdir fig_klempt2024_table4_blend.png}
  \caption[Test case 4.2, reproduction and paper]{Test case 4.2: vertical
    mid-plane, nutrient face at the bottom, reproduction next to Table~4 of
    Klempt et al.\ 2024 (CC BY 4.0), in the paper's colour bands.}
  \label{fig:app_klempt_table4}
\end{figure}

\section{Mechanics and units}

The reproduction has no mechanics. Two bounds show that the published curves
can nevertheless be compared with it. The material is incompressible, so the
current volume is $\alpha^3$ times the reference volume; with
$\alpha-1\le0.01$ at $T^*=1$ the volume-weighted means of $\phi$ and $c$
differ from the plain ones by at most $0.001$. The mechanical term of the
$\phi$ equation, which the paper drops, would be of order $10^4$ to $10^6$
per $T^*$ after division by the viscosity $\eta_\phi=10^{-10}$, against a
growth term of about $500$ per $T^*$; kept, it would stop the growth that
every figure shows.

The pressure plotted in Table~3 of the paper (of order $10^{-4}$~MPa) is
consistent with $\mu=3.3557$~MPa, as stated in the text, and not with the
$3.3557$~Pa of Table~2, which would give pressures a million times smaller.
This is an inference from the plot. The calculations of
Chapter~\ref{ch:ansys} use $E=10$~Pa as in Table~2; their stresses scale
linearly with $E$.

\section{Questions to the authors and their answers}

I sent these questions to the first author on 3 and 5 October 2026 and
received his answers on 5 October. They settle some points and leave others
open:
\begin{itemize}
  \item \emph{Front term.} His implementation uses Eq.~\eqref{eq:theory_phi}
        as printed: where the nutrient gradient behind the colony points
        towards it, the colony does not grow on that side. The isotropic
        share $w$ of the reproduction is therefore my change, needed to
        reach the published curves with the explicit scheme used here.
        Whether the finite element solution of the paper reaches them
        without it is open.
  \item \emph{Direction where the nutrient is uniform.} The implementation
        uses $\nabla_X c/(\|\nabla_X c\|+\varepsilon)$ with a small
        $\varepsilon$, of order $10^{-8}$ to $10^{-12}$.
  \item \emph{Consumption.} He recalls $g\,\phi$, as printed, but was not
        certain. The steady nutrient level of case 4.2 ``high'' (0.489 in
        the paper) follows from $g\,\phi\,c$ and not from $g\,\phi$
        (0.482 against 0.139, see above). This remains open until the
        original implementation is checked.
  \item \emph{Time.} Every simulation ran over a length of $1$. The time
        scales $s=5$ and $3$ that the reproduction needs for case 4.2 are
        therefore not explained by different reference times.
  \item \emph{Nutrient in case 4.1.} Only the corner, as assumed here.
  \item \emph{Diffusion coefficient $\beta$.} Chosen to give sensible
        results; it depends on the mesh. The value used in
        Chapter~\ref{ch:ansys}, Table~2 converted to the 2~mm cube, is
        therefore an assumption.
  \item \emph{Stiffness.} The unit in the paper may be wrong; he suggests
        $E=10$~kPa ($\mu\approx3.3$~kPa). This differs from the inference
        from the plot of Table~3 above ($\mu=3.3557$~MPa). The stresses of
        Chapter~\ref{ch:ansys} scale linearly with $E$, so either value
        rescales them without changing any comparison.
\end{itemize}
Whether $\alpha$ follows Eq.~36 or the update of Table~1
(Section~\ref{sec:theory_field}) remains to be asked. The first author has
offered his implementation; comparing it with the reproduction is part of
the continuation of this work.
